\SourcePage{232}{237}
\section{Nuclear radiation}

For nuclear $\gamma$ radiation, the dimensions of the system (the nuclear radius $R$) are generally small compared with the photon wavelength. The separations between nuclear levels, however, and hence the energy of a $\gamma$ quantum, are usually small compared with the energy per nucleon in the nucleus. The quantity $R/\lambda$ is therefore not directly related to the velocity $v/c$ of the nucleons in the nucleus and is, in general, much smaller. Correspondingly, the probability of $Ml$ radiation is usually greater than that of $E,l+1$ radiation (compare the beginning of \S~50).

The general selection rules for the total angular momentum (``spin'') of the nucleus and for parity are the same as for radiation by any system. A distinctive feature of nuclear radiation is the preva\SourcePageJoin{233}{238}lence of transitions of higher multipolarity. Unlike atoms, whose radiation is ordinarily electric dipole, nuclei at low energies make such transitions comparatively rarely because the selection rules forbid them.

If a radiative nuclear transition can be treated as a one-particle transition~--- a change in the state of one nucleon with the state of the nuclear ``core'' unchanged~--- additional selection rules apply to the angular momentum of that nucleon. Such ``one-particle'' selection rules, however, are obeyed with very low accuracy.

Selection rules involving isospin are specific to the nucleus. Recall that the projection $T_3$ of the isospin is already fixed by the mass number and nuclear charge:
\[
T_3=\frac12(Z-N)=Z-\frac A2.
\]
At a specified value of $T_3$, the magnitude of the isospin may take any value $T\geqslant|T_3|$. A selection rule for $T$ in radiative transitions arises because the operators of the electric and magnetic moments of the nucleus, when expressed through the nucleon isospin operators, are sums of a scalar and the $x_3$ component of a vector in isospin space (see III, \S~116). Their matrix elements are consequently nonzero only if
\begin{equation}
T'-T=0,\pm1.
\tag{55.1}
\end{equation}
By itself, however, this rule imposes no significant restriction on transitions in light nuclei (the only nuclei for which conservation of isospin is accurate enough to be meaningful): in fact, the low-lying levels of these nuclei include no levels with $T>1$.

For $E1$ transitions there is an additional rule, because the isoscalar part of the electric dipole moment drops out and its operator reduces to the $x_3$ component of an isospin vector (see III, \S~116). Thus, when $T_3=0$, transitions with $\Delta T=0$ are additionally forbidden. In other words, in nuclei having equal numbers of neutrons and protons ($N=Z$, $A=2Z$), $E1$ transitions are possible only when
\begin{equation}
T'-T=\pm1\qquad(T_3=0).
\tag{55.2}
\end{equation}
The accuracy of this rule naturally depends on how accurately the isospin of the nucleus is conserved.

The recoil of the nuclear core during the motion of individual nucleons also affects the probability of nuclear $E1$ transitions. As a result of this effect, protons contribute to the dipole moment with an effective charge $e(1-Z/A)$ rather than $e$, while neutrons contribute with a charge $-eZ/A$ rather than zero (see
\SourcePage{234}{239}
III, \S~118). The reduction of the proton's effective charge produces some suppression of the $E1$ transition probability.

The energy levels of nonspherical nuclei have rotational structure. Such nuclei therefore have a characteristic rotational structure in their $\gamma$-radiation spectra.

The symmetry of the field in which the nucleons move in a ``stationary'' nonspherical (axial) nucleus is the same as that of the field in which the electrons move in a ``stationary'' diatomic molecule composed of identical atoms (the point group $\mathrm C_{\infty h}$). The symmetry properties of the levels of a nonspherical nucleus, and with them the selection rules for matrix elements, are consequently analogous to those for the levels of a diatomic molecule (see III, \S~119). In particular, just as in a diatomic molecule made of identical atoms, electric-dipole transitions within the same rotational band (that is, without a change of the internal nuclear state) are forbidden; compare \S~54. Such transitions therefore occur as $E2$ or $M1$ transitions. In the first case the total nuclear angular momentum $J$ may change by 2 or 1, and in the second by 1.

According to (46.9), the quadrupole-transition probability summed over the values of the projection $M'$ of the total angular momentum of the nucleus in the final state is
\[
w_{E2}=\frac{\omega^5}{15\hbar c^5}\sum_{M'}
\left|\langle J'\Omega M'|Q_{2,-m}^{(\text{e})}|J\Omega M\rangle\right|^2
\]
($J$ is the total angular momentum of the nucleus, $\Omega$ its projection on the nuclear axis, and $m=M-M'$). Formula (110.8) (see III) expresses this sum in terms of the squares of specified quantities: the diagonal (in the internal nuclear state) transition quadrupole moments $\overline Q_{2\lambda}$, defined relative to the nuclear coordinate axes $\xi\eta\zeta$. Here $\lambda=\Omega-\Omega'$, so in the present case ($\Omega'=\Omega$) only the component $\overline Q_{20}$ occurs. The quantity called simply the quadrupole moment of the nucleus is defined by
\[
eQ_0=e\int\rho_{ii}(2\zeta^2-\xi^2-\eta^2)\,d\xi\,d\eta\,d\zeta
=-2e(\overline Q_{20})_{ii}.
\]
It follows that
\begin{equation}
w_{E2}(\Omega J\to\Omega J')
=\frac{\omega^5}{60\hbar c^5}Q_0^2(2j'+1)
\begin{pmatrix}J'&2&J\\-\Omega&0&\Omega\end{pmatrix}^{\!2}.
\tag{55.3}
\end{equation}
In expanded form,
\[
\begin{aligned}
w_{E2}(\Omega J\to\Omega,J-1)
&=\frac{\omega^5}{20\hbar c^5}Q_0^2
\frac{\Omega^2(J^2-\Omega^2)}
{(J-1)J(J+1)(2J+1)},\\
w_{E2}(\Omega J\to\Omega,J-2)
&=\frac{\omega^5}{40\hbar c^5}Q_0^2
\frac{(J^2-\Omega^2)[(J-1)^2-\Omega^2]}
{(J-1)J(2J-1)(2J+1)}.
\end{aligned}
\]

\SourcePage{235}{240}
The following qualification must, however, be made concerning these formulas. They use matrix elements calculated with wave functions of the form
\[
\psi_{J\Omega M}=\operatorname{const}\cdot\chi_\Omega D_{\Omega M}^{(J)}(\mathbf n)
\]
($\chi$ is the wave function of the internal nuclear state). These functions correspond to specified values, both in magnitude and sign, of the angular-momentum projection on the $\zeta$ axis. In a nucleus, however, the states have only a specified parity and magnitude of the angular-momentum projection (the latter is normally what is meant by $\Omega$). Thus, when $\Omega\ne0$, the initial and final wave functions should instead be combinations of the form
\[
\frac1{\sqrt2}(\psi_{J\Omega M}\pm\psi_{J,-\Omega,M}),\qquad
\frac1{\sqrt2}(\psi_{J'\Omega M'}\pm\psi_{J',-\Omega,M'}).
\]
The products of the first terms and of the second terms give the previous value of the quadrupole-moment matrix element. The ``cross'' products, on the other hand, give nonzero integrals when $2\Omega\leqslant l$\footnote{For matrix elements of $2^l$-pole moments, the integrands contain products of the form
\[
D_{-\Omega M}^{(J')*}D_{q'q}^{(l)}D_{\Omega M}^{(J)}.
\]
The angular integral is nonzero for $q'=-2\Omega$, whereas $q'$ ranges only from $-l$ to $l$; hence $2\Omega\leqslant l$.}. Strictly speaking, formula (55.3) is therefore inapplicable for $\Omega=1/2,1$; in these cases the transition probability contains an additional term that cannot be expressed through the mean quadrupole moment\footnote{In practice this term gives an appreciable correction only for $\Omega=1/2$, when the coupling between rotation and the internal state of the nucleus is especially strong (see III, \S~119).}.

By a derivation analogous to that of (55.3), the probability of an $M1$ transition is
\begin{equation}
\begin{aligned}
w_{M1}(\Omega J\to\Omega,J-1)
&=\frac{4\omega^3}{3\hbar c^3}\mu^2(2J-1)
\begin{pmatrix}J-1&1&J\\-\Omega&0&\Omega\end{pmatrix}^{\!2}\\
&=\frac{4\omega^3}{3\hbar c^3}\mu^2
\frac{J^2-\Omega^2}{J(2J+1)},
\end{aligned}
\tag{55.4}
\end{equation}
where $\mu$ is the magnetic moment of the nucleus (this formula is inapplicable for $\Omega=1/2$).
